\NormalFont\ShortTitle{1 Coríntios}
{\MT Primeira Carta de Paulo aos Coríntios

\par }\ChapOne{1}{\PP \VerseOne{1}Paulo, chamado apóstolo de Jesus Cristo pela vontade de Deus, e o irmão Sóstenes,
\VS{2}À igreja de Deus que está em Corinto, aos santificados em Cristo Jesus, chamados santos, com todos os que em todo lugar invocam o nome de nosso Senhor Jesus Cristo,
{\ADD{Senhor}} deles, e nosso;
\VS{3}Graça
{\ADD{seja}} convosco, e a paz de Deus nosso Pai, e do Senhor Jesus Cristo.
\VS{4}Sempre agradeço ao meu Deus por causa de vós, pela graça de Deus, que é dada a vós em Cristo Jesus.
\VS{5}Que em todas as coisas estais enriquecidos nele, em toda palavra, e em todo conhecimento;
\VS{6}Assim como o testemunho de Jesus Cristo foi confirmado entre vós.
\VS{7}De maneira que nenhum dom vos falta, enquanto esperais a manifestação do nosso Senhor Jesus Cristo,
\VS{8}o qual também vos firmará até o fim, irrepreensíveis no dia de nosso Senhor Jesus Cristo.
\VS{9}Fiel é Deus, por quem fostes chamados para a comunhão do seu Filho Jesus Cristo, nosso Senhor.
\VS{10}Mas eu vos rogo, irmãos, pelo nome do nosso Senhor Jesus Cristo, que todos faleis uma mesma coisa, e não haja divisões entre vós; antes estejais juntos com o mesmo entendimento, e com a mesma opinião.
\VS{11}Porque, meus irmãos, foi-me informado acerca de vós, pelos
{\ADD{da casa}} de Cloé, de que há brigas entre vós.
\VS{12}E digo isto, que cada um de vós diz: “Eu sou de Paulo” e “Eu
{\ADD{sou}} de Apolo” e “Eu
{\ADD{sou}} de Cefas” e “Eu
{\ADD{sou}} de Cristo”.
\VS{13}Cristo está dividido? Paulo foi crucificado por vós? Ou fostes vós batizados no nome de Paulo?
\VS{14}Agradeço a Deus que batizei nenhum de vós, a não ser Crispo e Gaio;
\VS{15}para que ninguém diga que eu tenha batizado em meu nome.
\VS{16}Também batizei aos da casa de Estefanas; além desses, não sei se batizei algum outro.
\VS{17}Porque Cristo não me enviou para batizar, mas sim, para evangelizar; não com sabedoria de palavras, para que a cruz de Cristo não se torne inútil.
\VS{18}Porque a palavra da cruz é loucura para os que perecem; mas para os que se salvam é poder de Deus.
\VS{19}Porque está escrito: Destruirei a sabedoria dos sábios, e aniquilarei a inteligência dos inteligentes.
\VS{20}Onde está o sábio? Onde está o escriba? Onde está o questionador desta era? Por acaso Deus não tornou a sabedoria deste mundo em loucura?
\VS{21}Pois já que, na sabedoria de Deus, o mundo não conheceu a Deus pela sabedoria, Deus se agradou de salvar os que creem por meio da loucura da pregação.
\VS{22}Porque os judeus pedem um sinal miraculoso, e os gregos buscam sabedoria.
\VS{23}Mas nós pregamos a Cristo crucificado,
{\ADD{que é}} motivo de ofensa para os Judeus, e loucura para os gregos.
\VS{24}Porém aos que são chamados, tanto judeus como gregos, Cristo
{\ADD{é}} poder de Deus e sabedoria de Deus.
\VS{25}Porque a loucura de Deus é mais sábia que os seres humanos; e a fraqueza de Deus é mais forte que os seres humanos.
\VS{26}Vede, pois, o vosso chamado, irmãos; pois dentre vós não
{\ADD{há}} muitos sábios em sabedoria humana,
\FTNT{*}{{\FR 1:26 }
sábios em sabedoria humana
Lit. sábios segundo a carne
} nem muitos poderosos, nem muitos da elite.
\VS{27}Mas Deus escolheu as coisas loucas deste mundo para envergonhar as sábias; e Deus escolheu as fracas deste mundo para envergonhar as fortes.
\VS{28}E Deus escolheu as coisas desprezíveis deste mundo, e as sem valor, e as que não são, para reduzir a nada as que são;
\VS{29}para que ninguém
\FTNT{†}{{\FR 1:29 }
Lit. nenhuma carne
} orgulhe de si mesmo diante dele.
\VS{30}Mas vós sois dele em Cristo Jesus, o qual, por parte de Deus, tornou-se para nós sabedoria, justiça, santificação, e resgate;
\VS{31}para que, assim como está escrito: “Aquele que se orgulha, orgulhe-se no Senhor”.

\par }\Chap{2}{\PP \VerseOne{1}E eu, irmãos, quando vim até vós, não vim vos anunciar o testemunho de Deus com excelência de palavras, ou sabedoria.
\VS{2}Porque não quis saber coisa alguma entre vós, a não ser Jesus Cristo, e ele crucificado.
\VS{3}E eu estive convosco em fraqueza, em temor, e em grande tremor.
\VS{4}E a minha palavra e a minha pregação não foram em palavras persuasivas de sabedoria humana, mas sim, em demonstração do Espírito e de poder;
\VS{5}para que a vossa fé não fosse em sabedoria humana, mas sim, no poder de Deus.
\VS{6}Contudo, falamos sabedoria entre os maduros, porém uma sabedoria não deste mundo, nem dos líderes deste mundo, que perecem.
\VS{7}Mas falamos a sabedoria de Deus, oculta em mistério, a qual Deus ordenou antes dos tempos para a nossa glória;
\VS{8}a
{\ADD{sabedoria}} que nenhum dos líderes deste mundo conheceu. Porque, se a tivessem conhecido, nunca teriam crucificado ao Senhor da glória.
\VS{9}Porém, assim como está escrito: “As coisas que o olho não viu, e não subiram que ao coração humano,
{\ADD{são}} as que Deus preparou para os que o amam”.
{\RQ{Isaías 64:4
}}
\VS{10}Mas Deus as revelou para nós pelo seu Espírito. Porque o Espírito investiga todas as coisas, até as profundezas de Deus.
\VS{11}Pois quem dentre os seres humanos sabe as coisas do ser humano, senão o espírito do ser humano que nele está? Assim também ninguém sabe as coisas de Deus, senão o Espírito de Deus.
\VS{12}Nós, porém, não recebemos o espírito do mundo, mas sim, o Espírito que provém de Deus; para que saibamos as coisas que por Deus nos são dadas;
\VS{13}as quais também falamos, não com palavras que a sabedoria ensina, mas com as que o Espírito Santo ensina, interpretando
{\ADD{coisas}} espirituais por
{\ADD{meios}} espirituais.
\VS{14}Mas o ser humano natural não compreende as coisas que são do Espírito de Deus, porque lhe são loucura; e não as pode entender, porque se compreendem espiritualmente.
\VS{15}Porém quem é espiritual compreende todas as coisas, mas não é compreendido por ninguém.
\VS{16}Pois quem conheceu a mente do Senhor, para que o possa instruir? Mas nós temos a mente de Cristo.

\par }\Chap{3}{\PP \VerseOne{1}E eu, irmãos, não pude vos falar como a espirituais, mas como a carnais, como a meninos em Cristo.
\VS{2}Com leite eu vos criei, e não com alimento sólido, porque não podíeis; nem mesmo agora ainda podeis;
\VS{3}porque ainda sois carnais. Pois, havendo entre vós inveja, brigas, e divisões, por acaso não sois carnais, e andais fazendo conforme o costume humano?
\VS{4}Porque dizendo um: “Eu sou de Paulo”, e outro: “Eu
{\ADD{sou}} de Apolo”, por acaso não sois carnais?
\VS{5}Ora, quem é Paulo, e quem é Apolo, senão servidores, pelos quais crestes, conforme o Senhor deu a cada um?
\VS{6}Eu plantei; Apolo regou; mas foi Deus quem deu o crescimento.
\VS{7}De maneira que nem o que planta é algo, nem o que rega; mas sim Deus, que dá o crescimento.
\VS{8}E o que planta e o que rega são um; mas cada um receberá a sua recompensa segundo o seu trabalho.
\VS{9}Porque somos cooperadores de Deus; vós sois a plantação de Deus, o edifício de Deus.
\VS{10}Segundo a graça de Deus que me foi dada, eu, como sábio construtor, pus o fundamento; e outro edifica sobre ele; mas cada um veja como edifica sobre ele.
\VS{11}Pois ninguém pode pôr fundamento diferente do que já está posto, o qual é Jesus Cristo.
\VS{12}E se alguém sobre este fundamento edificar ouro, prata, pedras preciosas, madeira, feno, palha,
\VS{13}a obra de cada um se manifestará; porque o dia a esclarecerá; pois pelo fogo se descobre; e qual é a obra de cada um, o fogo fará a prova.
\VS{14}Se a obra de alguém que construiu sobre ele permanecer, receberá recompensa.
\VS{15}Se a obra de alguém se queimar, ele sofrerá perda; porém o tal se salvará, todavia, como que
{\ADD{passado}} pelo fogo.
\VS{16}Não sabeis vós, que sois o templo de Deus? E que o Espírito de Deus habita em vós?
\VS{17}Se alguém destruir o templo de Deus, Deus ao tal destruirá; porque o templo de Deus é santo, o qual sois vós.
\VS{18}Ninguém se engane; se alguém entre vós nestes tempos pensa ser sábio, que se faça de louco, para que seja sábio.
\VS{19}Pois a sabedoria deste mundo é loucura para Deus; porque está escrito: Ele toma aos sábios pela sua própria astúcia.
\VS{20}E outra vez: O Senhor conhece os pensamentos dos sábios, que são vãos.
\VS{21}Portanto ninguém se orgulhe em seres humanos; porque tudo é vosso,
\VS{22}seja Paulo, seja Apolo, seja Cefas, seja o mundo, seja a vida, seja a morte, seja o presente, seja o futuro; tudo é vosso.
\VS{23}Mas vós sois de Cristo; e Cristo, de Deus.

\par }\Chap{4}{\PP \VerseOne{1}Que as pessoas nos considerem como servidores de Cristo, e guardiões dos mistérios de Deus.
\VS{2}Além disso, exige-se dos guardiões que cada um seja fiel.
\VS{3}Quanto a mim, pouco me importa ser julgado por vós, ou por qualquer juízo humano; nem eu julgo a mim mesmo.
\VS{4}Porque em nada me sinto culpável; mas nem por isso estou justificado; porém quem me julga é o Senhor.
\VS{5}Então nada julgueis antes do tempo, até que o Senhor venha, o qual também trará à luz as coisas ocultas nas trevas, e manifestará as intenções dos corações; e então cada um receberá de Deus a aprovação.
\VS{6}E isto, irmãos, fiz a comparação relativa a mim e a Apolo por causa de vós; para que em nós aprendais a não fazer suposições além do que está escrito; para que não fiqueis arrogantes por preferir um contra o outro.
\VS{7}Pois quem te faz diferentemente superior? E o que tens tu, que não tenhas recebido? E se o recebeste, por que te orgulhas, como se não o tivesse recebido?
\VS{8}Vós já estais satisfeitos, já estais ricos, e reinastes sem nós; e bom seria se vós reinásseis, para que também nós reinássemos convosco.
\VS{9}Pois acho que Deus colocou a nós, os apóstolos, por último, como sentenciados à morte; pois somos como espetáculo ao mundo, tanto aos anjos como aos seres humanos.
\VS{10}Nós
{\ADD{somos}} loucos por causa de Cristo, mas vós sábios em Cristo; nós fracos, mas vós fortes; vós ilustres, mas nós desonrados.
\VS{11}Até esta presente hora sofremos fome e sede, e estamos nus, e somos golpeados, e não temos morada certa.
\VS{12}E ficamos exaustos, trabalhando com as nossas próprias mãos; somos insultados, e bendizemos; somos perseguidos, e nós suportamos.
\VS{13}Somos blasfemados, e respondemos humildemente; somos feitos como a escória do mundo, e como o lixo de todos, até agora.
\VS{14}Não escrevo estas coisas para vos envergonhar; mas sim, para
{\ADD{vos}} advertir como a meus amados filhos.
\VS{15}Porque ainda que tivésseis dez mil tutores em Cristo, contudo não
{\ADD{tendes}} muitos pais. Porque eu vos gerei em Cristo Jesus por meio do Evangelho.
\VS{16}Portanto eu vos chamo para que sejais meus imitadores.
\VS{17}Por isso eu vos enviei Timóteo, que é meu filho amado e fiel no Senhor; ele vos lembrará dos meus caminhos em Cristo, como por todos os lugares eu ensino em cada igreja.
\VS{18}Mas alguns andam arrogantes, como se eu não fosse até vós.
\VS{19}Porém logo irei vos visitar, se o Senhor quiser; e
{\ADD{então}} conhecerei, não as palavras dos que andam arrogantes, mas sim o poder
{\ADD{deles}} .
\VS{20}Porque o Reino de Deus não
{\ADD{consiste}} de palavras, mas sim de poder.
\VS{21}Que quereis? Que eu venha até vós com a vara
{\ADD{da disciplina}} , ou com amor e espírito de mansidão?

\par }\Chap{5}{\PP \VerseOne{1}Totalmente se ouve
{\ADD{que há}} entre vós pecados sexuais, e tais pecados sexuais que nem mesmo entre os gentios se nomeia, como alguém que tomou como mulher a esposa de seu pai.
\VS{2}E vós estais orgulhosos! Ao invés disso, vós não deveríeis se entristecer, para que fosse tirado do meio de vós o que cometeu esta atitude?
\VS{3}Porque eu,
{\ADD{agindo}} como ausente de corpo, porém presente de espírito, já julguei como se estivesse presente, aquele que fez isto,
\VS{4}No nome do nosso Senhor Jesus Cristo, quando vós vos reunirdes, e o meu espírito junto de vós, com o poder de nosso Senhor Jesus Cristo,
\VS{5}De ao tal entregar a Satanás, para destruição da carne, para que o espírito seja salvo, no dia do Senhor Jesus.
\VS{6}Vosso orgulho não é bom. Não sabeis que um pouco de fermento faz levedar toda a massa?
\VS{7}Limpai pois o velho fermento, para que sejais nova massa, assim como vós sois
{\ADD{de fato}} sem fermento. Porque Cristo, nosso
{\ADD{cordeiro da}} páscoa, foi sacrificado por nós.
\VS{8}Portanto façamos festa, não no velho fermento, nem no fermento do mal e da perversão, mas em pães não levedados de sinceridade e de verdade.
\VS{9}Eu já vos escrevi por carta, que não vos mistureis com pecadores sexuais.
\VS{10}Porém não totalmente com os pecadores deste mundo, ou com os gananciosos, ou ladrões, ou
{\ADD{com os}} idólatras; porque então vos seria necessário sair do mundo.
\VS{11}Mas agora vos escrevo, que não vos mistureis se alguém, chamando a si de irmão, for pecador sexual, ou ganancioso, ou idólatra, ou maldizente, ou beberrão, ou ladrão. Com o tal nem mesmo comais.
\VS{12}Por que eu tenho que julgar também aos que estão de fora? Não julgais vós aos que estão de dentro?
\VS{13}Mas Deus julga aos que estão de fora. Tirai pois dentre vós a este maligno.

\par }\Chap{6}{\PP \VerseOne{1}Algum de vós, tendo um assunto contra outro, ousa ir a juízo perante os injustos, e não perante os santos?
\VS{2}Não sabeis que os santos julgarão ao mundo? E se o mundo é julgado por vós, por acaso sois vós indignos de julgar as coisas mínimas?
\VS{3}Não sabeis que julgaremos aos anjos? Quanto mais as coisas desta vida!
\VS{4}Pois se tiverdes assuntos em juízo, relativos a coisas desta vida, vós pondes como juízes aos que são menos estimados pela igreja?
\VS{5}Eu digo
{\ADD{isto}} para vos envergonhar. Não há entre vós algum sábio, nem pelo menos um, que não possa julgar entre seus irmãos?
\VS{6}Mas irmão vai a juízo contra outro irmão, e isto perante incrédulos?
\VS{7}Porém já é uma falta entre vós, por entre vós haverdes disputas. Por que não aceitais serdes injustiçados? Por que não aceitais o prejuízo?
\VS{8}Mas vós
{\ADD{mesmos}} injustiçais e prejudicais, e isto aos irmãos.
\VS{9}Ou não sabeis que os injustos não herdarão o Reino de Deus? Não erreis: nem os pecadores sexuais, nem os idólatras, nem os adúlteros, nem os sexualmente efeminados, nem os homens que fazem sexo com homens,
\VS{10}Nem os ladrões, nem os gananciosos, nem os beberrões, nem os maldizentes, nem os extorsores herdarão o Reino de Deus.
\VS{11}E isto alguns de vós éreis. Mas
{\ADD{já}} estais lavados, e santificados, e justificados no nome do Senhor Jesus, e pelo Espírito de nosso Deus.
\VS{12}Todas as coisas me são lícitas, mas nem todas as coisas convêm; todas as coisas me são lícitas, porém eu não deixarei me sujeitar por coisa alguma.
\VS{13}A comida é para o estômago, e o estômago é para a comida; mas Deus os aniquilará, tanto a um, como ao outro. Porém o corpo não é para o pecado sexual, mas para o Senhor, e o Senhor
{\ADD{é}} para o corpo.
\VS{14}E Deus ressuscitou ao Senhor, e também por seu poder nos ressuscitará.
\VS{15}Não sabeis vós que vossos corpos são membros de Cristo? Tomarei pois os membros de Cristo, e os farei membros de uma prostituta? De maneira nenhuma!
\VS{16}Ou não sabeis que o que se junta
{\ADD{sexualmente}} com uma prostituta se torna um
{\ADD{só}} corpo
{\ADD{com ela}} ? Porque
{\ADD{a Escritura}} diz: Os dois serão uma
{\ADD{só}} carne.
\VS{17}Mas o que se junta com o Senhor,
{\ADD{com ele}} se torna um
{\ADD{só}} espírito.
\VS{18}Fugi do pecado sexual. Todo pecado que o ser humano comete é fora do corpo; mas o que pratica pecado sexual, peca contra seu próprio corpo.
\VS{19}Ou não sabeis que vosso corpo é templo do Espírito Santo, o qual está em vós, o qual tendes de Deus, e
{\ADD{que}} não sois de vós mesmos?
\VS{20}Porque vós fostes comprados por alto preço; então glorificai a Deus em vosso corpo, e em vosso espírito, que são de Deus.

\par }\Chap{7}{\PP \VerseOne{1}E sobre as coisas que me escrevestes, bom é para o homem não tocar mulher.
\VS{2}Mas por causa dos pecados sexuais, tenha cada um sua própria mulher, e cada uma tenha seu próprio marido.
\VS{3}Que o marido satisfaça sua mulher como devido, e semelhantemente a mulher ao marido.
\VS{4}A mulher não tem poder sobre seu próprio corpo, mas sim o marido; e também da mesma maneira o marido não tem poder sobre seu próprio corpo, mas sim a mulher.
\VS{5}Não vos priveis um ao outro, a não ser por consentimento
{\ADD{de ambos}} por algum tempo, para que vos ocupeis com jejum, e para a oração; e voltai-vos outra vez a se juntarem, para que Satanás não vos tente, por causa de vossa falta de domínio próprio.
\VS{6}Mas isto vos digo como permissão, não como mandamento.
\VS{7}Porque eu queria que todos fossem como eu mesmo; mas cada um tem seu próprio dom de Deus, um de um jeito assim, e outro de um jeito diferente.
\VS{8}Mas digo aos solteiros, e às viúvas que lhes é bom se permanecessem como eu estou.
\VS{9}Mas se não podem se conter, casem-se; porque é melhor se casar do que se inflamar.
\VS{10}Porém aos casados eu mando, não eu, mas o Senhor, que a mulher não se separe do marido.
\VS{11}E se se separar, fique sem casar, ou se reconcilie com o marido; e que o marido não deixe a mulher.
\VS{12}Mas aos outros digo eu, não o Senhor: se algum irmão tem mulher descrente, e ela consente em habitar com ele, não a deixe.
\VS{13}E se alguma mulher tem marido descrente, e ele consente em habitar com ela, não o deixe.
\VS{14}Porque o marido descrente é santificado pela mulher; e a mulher descrente é santificada pelo marido. Pois de outra maneira vossos filhos seriam impuros; porém agora são santos.
\VS{15}Mas se o descrente se separar, separe-se
{\ADD{também}} . Em tal
{\ADD{situação}} o irmão ou a irmã não estão sujeitos à servidão; mas Deus nos chamou para a paz.
\VS{16}Porque o que tu sabes, mulher, se salvarás ao marido? Ou o que tu sabes, marido, se salvarás a mulher?
\VS{17}Porém cada um ande assim como Deus lhe repartiu, cada um como o Senhor o chamou; e assim ordeno em todas as igrejas.
\VS{18}É alguém chamado, estando já circunciso? Não desfaça a circuncisão. É alguém chamado estando ainda incircunciso? Não se circuncide.
\VS{19}A circuncisão nada vale, e a incircuncisão nada vale; mas
{\ADD{o que vale}} é a obediência aos mandamentos de Deus.
\VS{20}Cada um fique no propósito em que foi chamado.
\VS{21}Foste tu chamado sendo escravo? Não te preocupes com
{\ADD{isto}} ; mas se tu podes te tornares livre, aproveita.
\VS{22}Porque o que foi chamado no Senhor enquanto era escravo, para o Senhor é livre; da mesma maneira também, o que foi chamado, sendo livre, se torna escravo de Cristo.
\VS{23}Vós fostes comprados por alto preço;
{\ADD{portanto}} não vos façais escravos dos seres humanos.
\VS{24}Irmãos, cada um continue com Deus
{\ADD{no estado}} em que foi chamado.
\VS{25}E quanto às virgens, não tenho mandamento do Senhor; porém eu dou
{\ADD{minha}} opinião, dado que tenho recebido misericórdia do Senhor para que eu seja confiável.
\VS{26}Considero pois isto como bom, por causa da necessidade atual, que é bom à pessoa continuar assim como está.
\VS{27}Estás ligado à mulher, não busques se separar. Estás livre de mulher, não busques mulher.
\VS{28}Mas se também te casares, não pecas; e se a virgem se casar, não peca. Porém os tais terão aflições na carne, que eu
{\ADD{gostaria}} de livrá-los.
\VS{29}Porém digo isto, irmãos, porque o tempo que resta é breve; para que também os que tem mulheres, sejam como se não as tivessem;
\VS{30}E os que choram, como se não chorassem; e os que se alegram, como se não se alegrassem; e os que compram, como se não possuíssem.
\VS{31}E os que usam
{\ADD{das coisas}} deste mundo, como se
{\ADD{delas}} não abusassem; porque a aparência deste mundo é passageira.
\VS{32}E eu queria que estivésseis sem preocupações. O solteiro se preocupa com as coisas do Senhor, como irá agradar ao Senhor;
\VS{33}Porém o que é casado se preocupa com as coisas do mundo, como irá agradar à mulher.
\VS{34}A mulher casada e a virgem são diferentes: a que está para casar tem preocupação com as coisas do Senhor, para ser santa, tanto de corpo como de Espírito; mas a casada tem preocupação com as coisas do mundo, como irá agradar ao marido.
\VS{35}Porém digo isto para vosso próprio proveito; não para vos enlaçar, mas para
{\ADD{vos guiar}} ao que é decente, e servirem ao Senhor sem distração.
\VS{36}Mas se alguém lhe parece que indecentemente trata com sua virgem, se ela passar da idade da juventude, e assim convier se fazer, que tal faça o que quiser, não peca, casem-se.
\VS{37}Porém o que está firme em
{\ADD{seu}} coração, não tendo necessidade, mas tem poder sobre sua própria vontade, e isto decidiu em seu próprio coração, de guardar sua virgem, faz bem.
\VS{38}Portanto, o que
{\ADD{a dá}} em casamento, faz bem; mas o que não a dá em casamento, faz melhor.
\VS{39}A mulher casada está ligada pela Lei todo o tempo em que seu marido vive; mas se seu marido morrer, fica livre para se casar com quem quiser, contanto
{\ADD{que seja}} no Senhor.
\VS{40}Porém mais bem-aventurada é, se assim continuar, segundo minha opinião. E também eu penso, que tenho o Espírito de Deus.

\par }\Chap{8}{\PP \VerseOne{1}E quanto às coisas sacrificadas aos ídolos, sabemos que todos temos conhecimento. O conhecimento incha
{\ADD{de orgulho}} ,mas o amor edifica.
\VS{2}E se alguém acha saber alguma coisa, ainda nada conhece como se deve conhecer.
\VS{3}Mas se alguém ama a Deus, o tal dele é conhecido.
\VS{4}Portanto, quanto ao comer das coisas sacrificadas aos ídolos: sabemos que o ídolo nada é no mundo, e que não há qualquer outro Deus, a não ser um.
\VS{5}Porque, ainda que também haja
{\ADD{alguns}} que se chamem deuses, seja no céu, seja na terra (como há muitos deuses e muitos senhores),
\VS{6}Todavia para nós há
{\ADD{só}} um Deus, o Pai, do qual
{\ADD{são}} todas as coisas, e nós para ele; e um
{\ADD{só}} Senhor Jesus Cristo, pelo qual
{\ADD{são}} todas as coisas, e nós por ele.
\VS{7}Mas não há em todos
{\ADD{este}} conhecimento; porém alguns até agora comem com consciência do ídolo, como se
{\ADD{fossem}} sacrificadas aos ídolos; e sendo sua consciência fraca, fica contaminada.
\VS{8}Ora, a comida não nos faz agradáveis a Deus. Porque seja o que comamos, nada temos de mais; e seja o que não comamos, nada nos falta.
\VS{9}Mas tomai cuidado para que esta vossa liberdade não seja de maneira alguma escândalo para os fracos.
\VS{10}Porque se alguém vir a ti, que tens
{\ADD{este}} conhecimento, sentado
{\ADD{à mesa}} no templo dos ídolos, não será a consciência do que é fraco induzida a comer das coisas sacrificadas aos ídolos?
\VS{11}E por causa de teu conhecimento, perecerá assim o irmão fraco, pelo qual Cristo morreu?
\VS{12}Assim que, pecando contra os irmãos e ferindo sua fraca consciência, pecais contra Cristo.
\VS{13}Portanto, se a comida fizer cair ao meu irmão, nunca mais comerei carne, para que eu não faça ao meu irmão cair.

\par }\Chap{9}{\PP \VerseOne{1}Não sou eu apóstolo? Não sou livre? Não vi eu a Jesus Cristo nosso Senhor? Não sois vós minha obra no Senhor?
\VS{2}Se para os outros não sou apóstolo, ao menos para vós eu o sou; porque vós sois o selo de meu apostolado no Senhor.
\VS{3}Esta é a minha defesa para com os que me condenam.
\VS{4}Não temos nós poder de comer e de beber?
\VS{5}Não temos nós direito de trazer
{\ADD{conosco}} uma mulher irmã, como também os demais apóstolos, e os irmãos do Senhor, e Cefas?
\VS{6}Ou só eu, e Barnabé, não temos direito de não trabalhar?
\VS{7}Quem vai para a guerra a seu próprio custo? Quem planta a vinha, e não come do seu fruto? Ou quem apascenta o gado, e não come do leite do gado?
\VS{8}Digo eu isto segundo a lógica humana? Ou não diz a Lei também o mesmo?
\VS{9}Porque na Lei de Moisés está escrito: ao boi que trilha não atarás a boca. Por acaso Deus tem preocupação com os bois?
\VS{10}Ou totalmente diz por nós? Porque por nós isto está escrito; porque o que lavra, com esperança deve lavrar; e o que trilha com esperança, de sua esperança
{\ADD{deve}} ser participante.
\VS{11}Se nós vos semeamos as coisas espirituais, é muito que ceifemos as vossas
{\ADD{coisas}} carnais?
\VS{12}Se outros são participantes deste poder sobre vós,
{\ADD{porque}} não tanto mais nós? Mas nós não usamos deste poder; antes tudo suportamos, para não darmos impedimento algum ao Evangelho de Cristo.
\VS{13}Não sabeis vós, que os que administram as coisas sagradas, comem
{\ADD{daquilo que é}} sagrado? E os que continuamente estão junto ao altar, com o altar participam?
\VS{14}Assim também ordenou o Senhor, aos que anunciam o Evangelho, que vivam do Evangelho.
\VS{15}Porém eu de nenhuma destas coisas usei; e nem isto escrevi, para que assim se faça comigo; porque melhor me seria morrer, do que esvaecer este meu orgulho.
\VS{16}Porque se anunciar o Evangelho, para mim não é orgulho, pois a obrigação me é imposta. E ai de mim, se não anunciar o Evangelho!
\VS{17}Porque se eu faço de boa mente, tenho recompensa; mas se eu fizer de má vontade, a responsabilidade me é confiada.
\VS{18}Então que recompensa tenho? Que evangelizando, proponha o Evangelho de Cristo sem receber nada de volta, para não abusar do meu poder no Evangelho.
\VS{19}Porque, estando eu livre de todos, me fiz servo de todos, para ganhar
{\ADD{ainda}} mais.
\VS{20}E me fiz como judeu para os judeus, para ganhar aos judeus; como que estivesse debaixo da Lei, para os que estão debaixo da Lei, para ganhar aos que estão debaixo da Lei.
\VS{21}Ao que estão sem Lei, como se estivesse sem Lei (
{\ADD{porém}} não estando sem Lei para com Deus; mas para com Cristo debaixo da Lei), para ganhar os que estão sem Lei.
\VS{22}Me fiz de fraco para os fracos, para ganhar aos fracos; tudo me fiz para todos, para por todos os meios vir a salvar a alguns.
\VS{23}E isto eu faço por causa do Evangelho, para que eu também dele seja participante.
\VS{24}Não sabeis vós que os que correm nas competições, realmente todos correm, mas somente um leva o prêmio? Correi de tal maneira, que o alcanceis.
\VS{25}E todo aquele que luta, tenha domínio próprio sobre tudo. Pois aqueles fazem para receber uma coroa perecível, porém nós
{\ADD{para}} uma que não perece.
\VS{26}Então assim eu corro, não como
{\ADD{para um lugar}} incerto; assim luto, não como que dando socos no ar.
\VS{27}Em vez disso, eu subjugo meu corpo, e o reduzo à servidão, para que quando estiver pregando para os outros, eu mesmo não seja reprovado.

\par }\Chap{10}{\PP \VerseOne{1}Ora, irmãos, eu não quero que ignoreis que todos os nossos pais estiveram debaixo da nuvem, e todos pelo mar passaram;
\VS{2}E todos foram batizados na nuvem e no mar em Moisés;
\VS{3}E todos comeram de um mesmo alimento espiritual;
\VS{4}E todos beberam de uma mesma bebida espiritual; porque bebiam da pedra espiritual que
{\ADD{os}} seguia; e a pedra era Cristo.
\VS{5}Mas da maioria deles Deus não se agradou; por isso foram deixados no deserto.
\VS{6}E estas coisas nos servem de exemplos, para que não desejemos coisas ruins, como eles desejaram.
\VS{7}E não sejais idólatras, como alguns deles
{\ADD{foram}} ,como está escrito: O povo se sentou para comer, e para beber, e levantaram-se para se alegrarem.
\VS{8}E não pequemos sexualmente, como alguns deles assim pecaram, e em um dia vinte e três mil caíram.
\VS{9}E não tentemos a Cristo, como também alguns deles tentaram, e pereceram pelas serpentes.
\VS{10}E não murmureis, como também alguns deles murmuraram, e pereceram pelo destruidor.
\VS{11}E todas estas coisas lhes aconteceram como exemplos, e estão escritas para nosso aviso, em quem os fins dos tempos têm chegado.
\VS{12}Portanto, aquele que pensa estar de pé, olhe para que não caia.
\VS{13}Nenhuma tentação vos veio, que não fosse humana; porém Deus é fiel; que não vos deixará tentar mais do que o que podeis, antes com a tentação também dará a saída, para que a possais suportar.
\VS{14}Portanto, meus amados, fugi da idolatria.
\VS{15}Eu vos falo como que para prudentes; jugai vós
{\ADD{mesmos}} o que eu digo.
\VS{16}Por acaso o copo da bênção, que nós bendizemos, não é a comunhão do sangue de Cristo? O pão que partimos, não é a comunhão do corpo de Cristo?
\VS{17}Porque
{\ADD{assim como há}} um
{\ADD{só}} pão, nós muitos somos um
{\ADD{só}} corpo; porque todos participamos de um
{\ADD{só}} pão.
\VS{18}Vede a Israel segundo a carne: por acaso os que comem os sacrifícios não são participantes do altar?
\VS{19}Então o que eu digo? Que o ídolo é alguma coisa? Ou que o sacrifício ao ídolo seja alguma coisa?
\VS{20}Mas eu digo, que as coisas que os gentios sacrificam, são sacrificadas para os demônios, e não a Deus. E não quero que sejais participantes dos demônios.
\VS{21}Não podeis beber o copo do Senhor e o copo dos demônios; não podeis ser participantes da mesa do Senhor e da mesa dos demônios.
\VS{22}Por acaso
{\ADD{tentamos}} provocar ciúmes ao Senhor? Somos nós mais fortes do que ele?
\VS{23}Todas as coisas me são lícitas, mas nem todas as coisas convêm; todas as coisas me são licitas, mas nem todas as coisas edificam.
\VS{24}Ninguém busque para si próprio, antes cada um
{\ADD{busque o bem}} do outro.
\VS{25}De tudo o que se vende no açougue, comei, sem vos questionar por causa da consciência.
\VS{26}Porque a terra é do Senhor, e
{\ADD{também}} sua plenitude.
\VS{27}E se alguém dos descrentes vos convidar, e quiserdes ir, comei de tudo o que for posto diante de vós, sem vos questionar por causa da consciência.
\VS{28}Mas se alguém vos disser: “Isto é sacrifício a ídolos”, não comais, por causa daquele que
{\ADD{vos}} advertiu, e
{\ADD{por causa}} da consciência. Porque a terra e sua plenitude pertencem ao Senhor.
\VS{29}Mas digo da consciência do outro, não a tua. Por que, então, minha liberdade é julgada pela consciência do outro?
\VS{30}E se eu participo
{\ADD{da comida}} pela graça, por que sou ofendido naquilo que dou graças?
\VS{31}Portanto, ao comer, ou ao beber, ou ao fazer qualquer coisa, fazei tudo para a glória de Deus.
\VS{32}Sede sem escândalo, nem a judeus, nem a gregos, nem à Igreja de Deus.
\VS{33}Assim como eu também agrado a todos em tudo, não buscando meu próprio proveito, mas sim o de muitos, para que assim possam se salvar.

\par }\Chap{11}{\PP \VerseOne{1}Sede meus imitadores, assim como eu também sou de Cristo.
\VS{2}Eu vos louvo, irmãos, de que em tudo vos lembrais de mim, e retendes
{\ADD{minhas}} ordens, assim como eu vos entreguei.
\VS{3}Mas quero que saibais, que a cabeça de todo homem é Cristo; e a cabeça da mulher é o homem; e a cabeça de Cristo é Deus.
\VS{4}Todo homem que ora ou profetiza, tendo
{\ADD{alguma coisa}} sobre a cabeça, desonra sua própria cabeça.
\VS{5}Mas toda mulher que ora ou profetiza com a cabeça descoberta, desonra sua própria cabeça; porque é o mesmo que se a tivesse raspada.
\VS{6}Porque se a mulher não se cobre, raspe-se também; mas se para a mulher é vergonhoso se cortar ou se raspar, que se cubra.
\VS{7}Porque o homem não deve cobrir a cabeça, pois é a imagem e a glória de Deus; mas a mulher é a glória do homem.
\VS{8}Porque o homem não provem da mulher, mas sim a mulher do homem.
\VS{9}Porque também o homem não foi criado por causa da mulher, mas a mulher
{\ADD{foi criada}} por causa do homem.
\VS{10}Portanto a mulher deve ter sobre a cabeça sinal de autoridade, por causa dos anjos.
\VS{11}Porém, nem o homem é sem a mulher, nem a mulher sem o homem, no Senhor.
\VS{12}Porque assim como a mulher vem do homem, assim também é o homem pela mulher; porém tudo
{\ADD{vem}} de Deus.
\VS{13}Julgai-vos entre vós mesmos: é decente que a mulher ore a Deus descoberta?
\VS{14}Ou não vos ensina a mesma natureza, que ter cabelo longo é desonra para o homem?
\VS{15}Mas a mulher ter cabelo longo, lhe é honra, porque os cabelos lhe são dados por cobertura?
\VS{16}Porém se alguém parece ser inclinado a disputas, nós não temos tal costume, nem as igrejas de Deus.
\VS{17}Mas isto eu vos declaro, que não vos louvo, porque vós vos ajuntais, não para melhor, mas para pior.
\VS{18}Porque, em primeiro lugar, eu ouço que, quando vos ajuntais, há divisões entre vós; e em parte eu acredito.
\VS{19}Porque é necessário que haja diferenças de opinião
\FTNT{*}{{\FR 11:19 }
diferenças de opinião
ou: facções. Tradicionalmente: heresias
} entre vós, para que os que são aprovados sejam evidenciados entre vós.
\VS{20}Quando, pois, vós vos reunis, não é para comer da ceia do Senhor;
\VS{21}Porque cada um, comendo antes, toma sua própria ceia; e um fica com fome, e outro fica bêbado.
\VS{22}Por acaso não tendes casas para comer e para beber? Ou desprezais a Igreja de Deus, e envergonhais os que não têm? Que vos direi? Eu vos elogiarei? Nisto eu não
{\ADD{vos}} elogio.
\VS{23}Porque eu recebi do Senhor o que também vos entreguei; que o Senhor Jesus, na noite em que foi traído, tomou o pão;
\VS{24}E tendo dado graças, o partiu, e disse: Tomai, comei; isto é o meu corpo, que é partido por vós; fazei isto em memória de mim.
\VS{25}Semelhantemente também, depois de cear,
{\ADD{tomou}} o copo, dizendo: Este copo é o novo Testamento em meu sangue. Fazei isto todas as vezes que
{\ADD{o}} beberdes, em memória de mim.
\VS{26}Porque todas as vezes que comerdes deste pão, e beberdes deste copo, anunciai a morte do Senhor, até que ele venha.
\VS{27}Portanto quem comer deste pão, ou beber deste copo do Senhor indignamente, será culpado do corpo e do sangue do Senhor.
\VS{28}Portanto examine-se cada um a si mesmo, e assim coma deste pão, e beba deste copo.
\VS{29}Porque quem come e bebe indignamente, para si mesmo come e bebe condenação, não discernindo o corpo do Senhor.
\VS{30}Por esta causa há entre vós muitos fracos e doentes, e muitos que dormem.
\FTNT{†}{{\FR 11:30 }
Isto é, muitos que morrem
}
\VS{31}Porque se nós julgássemos a nós mesmos, não seríamos julgados.
\VS{32}Mas quando somos julgados, somos repreendidos pelo Senhor; para que não sejamos condenados com o mundo.
\VS{33}Portanto, meus irmãos, quando vos ajuntardes para comer, esperai uns aos outros.
\VS{34}Porém se alguém tiver fome, coma em casa; para que não vos ajunteis para
{\ADD{sofrerdes}} julgamento. Quanto às demais coisas, as ordenarei quando vier.

\par }\Chap{12}{\PP \VerseOne{1}E quanto aos
{\ADD{dons}} espirituais, irmãos, não quero que sejais ignorantes.
\VS{2}Vós sabeis que éreis gentios, vós vos desviáveis para os ídolos mudos, conforme éreis guiados.
\VS{3}Por isso eu vos deixo claro que ninguém que fala pelo Espírito de Deus chama a Jesus de maldito;
\FTNT{*}{{\FR 12:3 }
equiv. anátema
} e ninguém pode dizer
{\ADD{que}} Jesus
{\ADD{é}} o Senhor, a não ser pelo Espírito Santo.
\VS{4}E há variedade de dons, mas o Espírito é o mesmo.
\VS{5}E há variedade de trabalhos
{\ADD{na Igreja}} ,mas o Senhor é o mesmo.
\VS{6}E há variedade de operações, mas é o mesmo Deus que opera tudo em todos.
\VS{7}Mas a cada um é dada a manifestação do Espírito, para o que for conveniente.
\VS{8}Porque a um é dada, pelo Espírito, palavra de sabedoria; e a outro palavra de conhecimento, pelo mesmo Espírito.
\VS{9}e a outro fé pelo mesmo Espírito; e a outro dons de curas, pelo mesmo Espírito.
\VS{10}E a outro operações de milagres; e a outro profecia; e a outro o
{\ADD{dom de}} discernir aos espíritos; e a outro variedade de línguas; e a outro interpretação de línguas.
\VS{11}Mas todas estas coisas
{\ADD{quem}} opera
{\ADD{é}} um e o mesmo Espírito, repartindo particularmente a cada um como quer.
\VS{12}Porque assim como o corpo é um, e tem muitos membros, e todos os membros deste único corpo, sendo muitos, compõem um só corpo; assim também é Cristo.
\VS{13}Porque também todos nós somos batizados em um
{\ADD{só}} Espírito, para um
{\ADD{só}} corpo, quer judeus, quer gregos, quer servos, quer livres; e para todos nós foi dado de beber um
{\ADD{só}} Espírito.
\VS{14}Porque também o corpo não é um
{\ADD{só}} membro, mas muitos.
\VS{15}Se o pé disser: Por eu não ser mão, não sou do corpo;Por isso, ele não é do corpo?
\VS{16}E se a orelha disser: Por eu não ser olho, não sou do corpo;Por isso, ela não é do corpo?
\VS{17}Se todo o corpo fosse olho, onde
{\ADD{estaria}} o ouvido? Se
{\ADD{fosse}} todo ouvido, onde
{\ADD{estaria}} o olfato?
\VS{18}Mas agora Deus pôs os membros no corpo, a cada qual deles como ele quis.
\VS{19}Porque se todos fossem um
{\ADD{só}} membro, onde
{\ADD{estaria}} o corpo?
\VS{20}Mas há agora muitos membros, porém um
{\ADD{só}} corpo.
\VS{21}E o olho não pode dizer à mão: Não tenho necessidade de ti;Nem a cabeça aos pés: Não tenho necessidade de vós;
\VS{22}Antes até os membros do corpo que
{\ADD{nos}} parecem ser os mais fracos, são muito mais necessários.
\VS{23}E os que pensamos serem os menos honrados do corpo, a esses muito mais honramos; e os nossos mais desprezíveis têm muito mais respeito.
\VS{24}Porém os nosso mais respeitáveis não têm
{\ADD{tanta}} necessidade; mas
{\ADD{assim}} Deus juntou o corpo, dando muito mais honra ao que tinha falta
{\ADD{dela}} .
\VS{25}Para que não haja divisão no corpo, mas que os membros tenham igual cuidado uns com os outros.
\VS{26}E seja que,
{\ADD{caso}} um membro sofra,
{\ADD{também}} os membros sofrem juntos; e
{\ADD{caso}} um membro seja honrado, todos os membros se alegram juntos.
\VS{27}E vós sois o corpo de Cristo, e membros em particular.
\VS{28}E Deus pôs a uns na Igreja: em primeiro lugar apóstolos, em segundo profetas, em terceiro mestres; depois milagres, dons de curas, socorros, lideranças, variedades de línguas.
\VS{29}Por acaso são todos apóstolos? São todos profetas? São todos mestres? São todos com dons de milagres?
\VS{30}Tem todos dons de curas? Falam todos
{\ADD{várias}} línguas? Interpretam todos?
\VS{31}Porém desejai com zelo pelos melhores dons; e eu vos mostro um caminho ainda mais excelente.

\par }\Chap{13}{\PP \VerseOne{1}Ainda que eu falasse as línguas dos seres humanos e dos anjos, e não tivesse amor, seria como o metal que soa, ou como o sino que retine.
\VS{2}E ainda que tivesse
{\ADD{o dom}} de profecia, e soubesse todos os mistérios, e todo o conhecimento; e ainda que tivesse toda a fé, de tal maneira que movesse os montes de lugar, e não tivesse amor, nada seria.
\VS{3}E ainda que eu distribuísse todos os meus bens para alimentar
{\ADD{aos pobres}} , e ainda que entregasse meu corpo para ser queimado, e não tivesse amor, nada me aproveitaria.
\VS{4}O amor é paciente, é bondoso; o amor não é invejoso; o amor não é orgulhoso, não é arrogante.
\VS{5}O amor não trata mal; não busca os próprios
{\ADD{interesses}} , não se ira, não é rancoroso.
\VS{6}Não se alegra com a injustiça, mas se alegra com a verdade.
\VS{7}Tudo aguenta, tudo crê, tudo espera, tudo suporta.
\VS{8}O amor nunca falha. Porém as profecias serão aniquiladas; as línguas acabarão, e o conhecimento será aniquilado.
\VS{9}Porque em parte conhecemos, e em parte profetizamos;
\VS{10}Mas quando vier o
{\ADD{que é}} completo,{\ADD{Ou: perfeito}} então o que é em parte será aniquilado.
\VS{11}Quando eu era menino, falava como menino, sentia como menino, pensava como menino; mas quando me tornei homem, aniquilei as coisas de menino.
\VS{12}Porque agora vemos por espelho em enigma, mas então
{\ADD{veremos}} face a face; agora conheço em parte, mas então conhecerei assim como sou conhecido.
\VS{13}E agora continuam a fé, a esperança,
{\ADD{e}} o amor, estes três; porém o maior destes é o amor.

\par }\Chap{14}{\PP \VerseOne{1}Segui o amor, e desejai com zelo pelos
{\ADD{dons}} espirituais; porém principalmente que profetizeis.
\VS{2}Porque o que fala língua
{\ADD{estranha}} , não fala aos seres humanos, mas sim a Deus; porque ninguém
{\ADD{o}} entende, mas em espírito fala mistérios.
\VS{3}Mas o que profetiza, fala aos seres humanos
{\ADD{para}} edificação, convencimento, e consolação.
\VS{4}O que fala língua
{\ADD{estranha}} edifica a si mesmo; mas o que profetiza edifica à igreja.
\VS{5}E eu quero que todos vós faleis línguas, porém mais
{\ADD{ainda}} que profetizeis; porque o que profetiza é maior que o que fala línguas, a não ser que
{\ADD{também}} interprete, para que a igreja receba edificação.
\VS{6}E agora irmãos, se eu viesse até vos falando línguas, o que vos aproveitaria, se não vos falasse ou por revelação, ou por conhecimento, ou por profecia, ou por doutrina?
\VS{7}E até as coisas inanimadas, que produzem som, seja flauta, seja harpa, se não derem distinção de sons, como saberá o que se toca com a flauta, ou com a harpa?
\VS{8}Porque também se a trombeta der som incerto, quem se preparará para a guerra?
\VS{9}Assim
{\ADD{mesmo}} também vós, se com a língua não derdes palavra compreensível, como se entenderá o que se diz? Porque estareis falando para o ar.
\VS{10}Por exemplo, há tantos tipos de vozes no mundo, e nenhuma delas é muda.
\VS{11}Portanto, se eu não souber o propósito da voz, serei estrangeiro para o que fala, e o que fala será estrangeiro para mim.
\VS{12}Assim também vós, dado que desejais os dons espirituais, procurai
{\ADD{neles}} abundar, para edificação da igreja.
\VS{13}Portanto, o que fala em língua
{\ADD{estranha}} , ore para que possa interpretar.
\VS{14}Porque se eu orar em língua
{\ADD{estranha}} , meu espírito ora, mas meu entendimento fica sem fruto.
\VS{15}Então é o que? Orarei com o espírito, mas também orarei com o entendimento; cantarei com o espírito, mas também cantarei com o entendimento.
\VS{16}De outra maneira, se tu bendisseres com o espírito, como aquele que não tem conhecimento dirá amém por teu bendizer? Pois ele não sabe o que tu dizes.
\VS{17}Porque em verdade tu bem dás graças; mas o outro não é edificado.
\VS{18}Graças dou a meu Deus, que mais línguas falo que todos vós.
\VS{19}Porém eu quero
{\ADD{mais}} falar na igreja cinco palavras com meu entendimento, para que eu também possa instruir aos outros, do que dez mil palavras em língua
{\ADD{estranha}} .
\VS{20}Irmãos, não sejais meninos no entendimento, mas sede meninos na malícia, e adultos no entendimento.
\VS{21}Na Lei está escrito: A este povo falarei por gente de outras línguas, e por outros lábios; e ainda assim não me ouvirão, diz o Senhor.
\VS{22}Então as línguas
{\ADD{estranhas}} são por sinal, não para os que creem, mas para os que não creem; e a profecia não é para os que não creem, mas para os que creem.
\VS{23}Pois se toda a igreja se reunir, e todos falarem línguas
{\ADD{estranhas}} , e entrarem alguns sem entendimento ou descrentes, não dirão que estais loucos?
\VS{24}Mas se todos profetizarem, e algum sem entendimento ou descrente entrar, por todos é convencido,
{\ADD{e}} por todos é julgado.
\VS{25}E assim os segredos de seu coração ficam manifestos, e assim, lançando-se sobre
{\ADD{seu}} rosto, adorará a Deus, reconhecendo publicamente que Deus está entre vós.
\VS{26}Então o que há, irmãos? Quando vos reunis, tem cada um de vós salmo, tem doutrina, tem língua
{\ADD{estranha}} , tem revelação, tem interpretação, tudo se faça para edificação.
\VS{27}E se alguém falar em língua
{\ADD{estranha}} ,
{\ADD{sejam}} dois, ou no máximo três, e alternando-se, e um que interprete.
\VS{28}Mas se não houver intérprete, cale-se na igreja, fale porém consigo mesmo, e com Deus.
\VS{29}E falem dois ou três profetas, e os outros julguem.
\VS{30}Mas se a outro, que estiver sentada, for revelada
{\ADD{alguma coisa}} , cale-se o primeiro.
\VS{31}Porque todos vós podeis profetizar, um após o outro, para que todos aprendam, e sejam todos fortalecidos.
\VS{32}E os espíritos dos profetas estão sujeitos aos profetas.
\VS{33}Porque Deus não é de confusão, mas de paz, como em todas as igrejas dos santos.
\VS{34}Vossas mulheres fiquem caladas nas igrejas; porque não lhes é permitido falar, mas
{\ADD{que}} estejam sujeitas, como também a Lei o diz.
\VS{35}E se quiserem aprender alguma coisa, perguntem a seus próprios maridos em casa; porque é impróprio as mulheres falarem na igreja.
\VS{36}Por acaso a palavra de Deus saiu de vós? Ou ela somente chegou a vós?
\VS{37}Se alguém pensa ser profeta, ou espiritual, reconheça que as coisas que escrevo são mandamento do Senhor.
\VS{38}Porém se alguém ignora, que ignore.
\VS{39}Portanto, irmãos, desejai com zelo profetizar, e não impeçais o falar línguas
{\ADD{estranhas}} .
\VS{40}Fazei tudo decentemente e com ordem.

\par }\Chap{15}{\PP \VerseOne{1}Também irmãos, eu vos declaro o Evangelho, que já vos tenho anunciado, o qual também recebestes, no qual também estais.
\VS{2}Pelo qual também sois salvos, se retiverdes a palavra naquela maneira, em que eu vos tenho anunciado; a não ser se tenhais crido em vão.
\VS{3}Porque primeiramente vos entreguei o que também recebi, que Cristo morreu por nossos pecados, segundo as Escrituras;
\VS{4}E que foi sepultado, e que ressuscitou ao terceiro dia, segundo as Escrituras;
\VS{5}E que foi visto por Cefas, depois pelos doze.
\VS{6}Depois foi visto de uma fez por mais de quinhentos irmãos, dos quais a maioria
{\ADD{ainda}} vive, e também alguns já dormem.
\FTNT{*}{{\FR 15:6 }
Isto é, alguns já morreram
}
\VS{7}Depois foi visto por Tiago; depois por todos os apóstolos.
\VS{8}E por último de todos também foi visto por mim, como que a um nascido fora do tempo.
\VS{9}Porque eu sou o menor dos apóstolos, que não sou digno de ser chamado apóstolo, porque persegui a Igreja de Deus.
\VS{10}Mas pela graça de Deus sou o que sou; e sua graça
{\ADD{concedida}} a mim não foi vã; ao invés disso trabalhei muito mais que todos eles; porém não eu, mas a graça de Deus que está comigo.
\VS{11}Portanto, seja eu, sejam eles, assim pregamos, e assim crestes.
\VS{12}E se é pregado que Cristo ressuscitou dos mortos, como dizem alguns dentre vós, que não há ressurreição dos mortos?
\VS{13}E se não há ressurreição dos mortos,
{\ADD{então}} Cristo não ressuscitou.
\VS{14}E se Cristo não ressuscitou, então é vã nossa pregação, e vã também é vossa fé.
\VS{15}E assim também somos encontrados como falsas testemunhas de Deus; pois testemunhamos de Deus, que ressuscitou a Cristo, ao qual
{\ADD{porém}} não ressuscitou, se na verdade os mortos não ressuscitam.
\VS{16}Porque se os mortos não ressuscitam,
{\ADD{então}} Cristo também não ressuscitou.
\VS{17}E se Cristo não ressuscitou, vossa fé é vá,
{\ADD{e}} ainda estais em vossos pecados.
\VS{18}Portanto também pereceram os que dormiram em Cristo.
\VS{19}Se nesta vida esperamos somente em Cristo, somos os mais miseráveis de todas as pessoas.
\VS{20}Mas de fato Cristo ressuscitou dos mortos,
{\ADD{e}} foi feito as primícias dos que dormiram.
\VS{21}Pois dado que a morte
{\ADD{veio}} por um homem, também por um homem
{\ADD{veio}} a ressurreição dos mortos.
\VS{22}Porque assim como em Adão todos morrem, assim também em Cristo todos serão vivificados.
\VS{23}Mas cada um em sua ordem; Cristo, as primícias; depois os que são de Cristo, em sua vinda.
\VS{24}Depois será o fim, quando entregar o Reino de Deus e ao Pai, quando aniquilar todo domínio, e toda autoridade e poder.
\VS{25}Porque convém que ele reine, até que tenha posto todos os inimigos debaixo de seus pés.
\VS{26}O último inimigo, que será aniquilado,
{\ADD{é}} a morte.
\VS{27}Porque ele sujeitou todas as coisas debaixo de seus pés. Porém quando diz que todas as coisas
{\ADD{lhe}} estão sujeitas, está claro que se excetua aquele que todas as coisas lhe sujeitou.
\VS{28}E quando todas as coisas lhe forem sujeitas, então também o mesmo Filho se sujeitará a aquele, que todas as coisas lhe sujeitou, para que Deus seja tudo em todos.
\VS{29}De outra maneira, que farão os que se batizam pelos mortos, se totalmente os mortos não ressuscitam? Por que, então, se batizam pelos mortos?
\VS{30}Por que também nós à toda hora estamos em perigo?
\VS{31}A cada dia eu morro pelo orgulho que tenho por vós, em Cristo Jesus nosso Senhor.
\VS{32}Se por motivo humano lutei contra feras em Éfeso, qual é o meu proveito? Se os mortos não ressuscitam, comamos e bebamos, que amanhã morreremos.
\VS{33}Não erreis. As más companhias corrompem os bons costumes.
\VS{34}Despertai para a justiça, e não pequeis; porque alguns
{\ADD{ainda}} não têm o conhecimento de Deus. Eu digo isto para vossa vergonha.
\VS{35}Mas alguém dirá: Como ressuscitarão os mortos? E com que corpo virão?
\VS{36}Louco, o que tu semeias não é vivificado, se
{\ADD{primeiro}} não morrer.
\VS{37}E o que semeias, não semeias o corpo que irá brotar, mas o grão nu, como o de trigo, ou de qualquer outro
{\ADD{grão}} .
\VS{38}Mas Deus lhe dá o corpo como quer, e a cada semente seu próprio corpo.
\VS{39}Toda carne não é a mesma carne; mas uma é a carne humana, e outra a carne dos animais, e outra a dos peixes, e outra a das aves.
\VS{40}E há corpos celestiais, e corpos terrestres; mas uma é a glória dos celestiais, e outra a dos terrestres.
\VS{41}Outra é a glória do Sol, e outra a glória da Lua, e outra a glória das estrelas; porque
{\ADD{uma}} estrela difere em glória
{\ADD{de outra}} estrela.
\VS{42}Assim também será a ressurreição dos mortos. Semeia-se
{\ADD{o corpo}} degradável, ressuscitará capaz de nunca se degradar.
\VS{43}Semeia-se em desonra, ressuscitará em glória. Semeia-se em fraqueza, ressuscitará em força.
\VS{44}Semeia-se corpo natural, ressuscitará corpo espiritual. Há corpo animal, e há corpo espiritual.
\VS{45}Assim também está escrito: O primeiro homem, Adão, foi feito em alma vivente; o último Adão em espírito vivificante.
\VS{46}Mas não é primeiro o espiritual; mas sim o carnal, depois o espiritual.
\VS{47}O primeiro homem da terra
{\ADD{é}} terrestre; o segundo homem,
{\ADD{que é}} o Senhor,
{\ADD{é}} do Céu.
\VS{48}Assim como
{\ADD{é}} o terrestre, também
{\ADD{são}} os terrestres; e assim como
{\ADD{é}} o celestial, também
{\ADD{são}} os celestiais.
\VS{49}E assim como trouxemos a imagem do terrestre, assim também traremos a imagem do celestial.
\VS{50}Porém isto digo, irmãos, que a carne e o sangue não podem herdar o Reino de Deus, nem o que é degradável herda a capacidade de não se degradar.
\VS{51}eis que eu vos digo um mistério: nem todos em verdade dormiremos; porém todos seremos transformados.
\VS{52}Em um momento, em um piscar de olhos, à ultima trombeta; porque a trombeta soará, e os mortos ressuscitarão capazes de não se degradar, e nós seremos transformados.
\VS{53}Porque é necessário que este
{\ADD{corpo}} degradável revista-se da capacidade de não se degradar, e este
{\ADD{corpo}} mortal vista a imortalidade.
\VS{54}E quando este
{\ADD{corpo}} se revestir da capacidade de não se degradar, e este
{\ADD{corpo}} mortal se vestir de imortalidade, então se cumprirá a palavra que está escrita: Na vitória, a morte é tragada.
\VS{55}Onde está, ó morte, o teu aguilhão? Onde está, ó Xeol,
\FTNT{†}{{\FR 15:55 }
Xeol é o lugar dos mortos
} a tua vitória?
\VS{56}O aguilhão da morte é o pecado, e a força do pecado é a Lei.
\VS{57}Mas graças a Deus, que nos dá vitória por meio do nosso Senhor Jesus Cristo.
\VS{58}Portanto, meus amados irmãos, sede firmes, imóveis
{\ADD{e}} sempre abundantes na obra do Senhor, sabendo que o vosso trabalho não é vão no Senhor.

\par }\Chap{16}{\PP \VerseOne{1}E quanto à coleta que
{\ADD{se faz}} para os santos, fazei também como ordenei às igrejas da Galácia.
\VS{2}Cada primeiro
{\ADD{dia}} da semana, cada um de vós ponha
{\ADD{alguma coisa}} à pate, reservando
{\ADD{para isso}} conforme a prosperidade que tiver obtido, para que, quando eu vier, não se façam coletas.
\VS{3}E quando eu vier, enviarei aos que por cartas aprovardes, para que levem vossa doação para Jerusalém.
\VS{4}E se for necessário que eu também vá, irão comigo.
\VS{5}Mas virei até vós, depois de passar pela Macedônia (porque passarei pela Macedônia).
\VS{6}E pode ser que eu fique
{\ADD{um tempo}} convosco ou também passar o inverno, para que me envieis para onde quer que eu for.
\VS{7}Porque eu não quero vos ver agora
{\ADD{apenas}} de passagem; mas espero ficar convosco por algum tempo, se o Senhor o permitir.
\VS{8}Porém ficarei em Éfeso até o Pentecostes.
\VS{9}Porque uma porta grande e eficaz se abriu para mim, e há muitos adversários.
\VS{10}E se vier Timóteo, olhai para que ele esteja sem temor convosco; porque assim como eu, ele também faz a obra do Senhor.
\VS{11}Portanto ninguém o despreze; mas enviai-o em paz, para que ele venha a mim; porque eu o espero com os irmãos.
\VS{12}E quanto ao irmão Apolo, muito roguei que ele viesse com os irmãos até vós; mas ele de maneira nenhuma teve vontade de vir por agora; porém oferecendo-se a ele boa oportunidade, ele virá.
\VS{13}Vigiai, ficai firmes na fé, sede corajosos,
{\ADD{e}} vos esforçai.
\VS{14}Todas as vossas coisas sejam feitas em amor.
\VS{15}Vós sabeis que
{\ADD{os da}} casa de Estéfanas foram os primeiros da Acaia, e que eles têm se dedicado ao trabalho dos santos; rogo-vos, pois, irmãos,
\VS{16}Que também vos sujeiteis a eles, e a todo aquele que opera e trabalha junto
{\ADD{conosco}} .
\VS{17}E eu me alegro com a vinda de Estéfanas, e de Fortunato, e de Acaico, porque estes supriram o que de vossa
{\ADD{parte me}} faltava.
\VS{18}Porque eles confortaram meu espírito, e
{\ADD{também}} o vosso. Então reconhecei aos tais.
\VS{19}As igrejas da Ásia vos saúdam. Áquila e Priscila vos saúdam-vos afetuosamente no Senhor, e também a igreja que está na casa deles.
\VS{20}Todos os irmãos vos saúdam. Saudai-vos uns as outros com beijo santo.
\VS{21}Saudação de minha
{\ADD{própria}} mão, de Paulo.
\VS{22}Se alguém não ama ao Senhor Jesus Cristo seja anátema. Vem, Senhor.
\VS{23}A graça do Senhor Jesus Cristo esteja convosco.
\VS{24}Meu amor esteja com todos vós em Cristo Jesus. Amém.
{\ADD{A primeira Carta aos Coríntios foi escrita de Filipos, e enviada por Estéfanas, Fortunato, Acaico e Timóteo}}
\par }